Ce chapitre est le seul du livre à ne pas parler de mathématiques, et le seul à n'être
qu'entamé. Il prolonge l'algorithmique du collège dans la spécialité \emph{numérique et
sciences informatiques} du lycée, dont les quarante et un énoncés sont posés dans l'index
du dépôt ; huit sont démontrés ici, ceux qui portent sur la représentation des données.

Le choix de commencer par là n'est pas un hasard. Un ordinateur ne manipule que des suites
de bits : tout le reste \textemdash{} les entiers, les nombres à virgule, les caractères
\textemdash{} est convention. Ces conventions se démontrent, et c'est même la partie de la
spécialité qui se formalise le plus directement, parce qu'elle ne parle que d'entiers.
L'écriture binaire d'un entier existe, elle est unique, elle se relit comme une somme de
puissances de deux : trois énoncés qu'un cours présente comme évidents et qui demandent
chacun une démonstration.

Deux résultats méritent qu'on s'y arrête. Le complément à deux explique pourquoi ajouter
$1$ au plus grand entier positif d'un mot machine rend le plus petit négatif, sans erreur
ni avertissement : le calcul est juste modulo $2^n$, et faux dans $\mathbb{Z}$. Et
l'absence d'écriture binaire finie pour un dixième, démontrée ici comme on démontre
l'irrationalité de $\sqrt2$, est la raison pour laquelle $0{,}1 + 0{,}2$ ne rend pas
exactement $0{,}3$ \textemdash{} la première surprise de qui commence à programmer.

Ce qui manque est d'une autre nature que dans les chapitres de mathématiques. Les tris, la
recherche dichotomique, les arbres, les graphes ne demandent pas une théorie hors de
portée : ils demandent d'écrire les algorithmes eux-mêmes, puis un modèle de coût pour
énoncer ce qu'ils coûtent. C'est un travail de programmation avant d'être un travail de
démonstration, et il reste à faire.
